% !TEX root = ../frode-bibelen.tex
% ----------------------------------------------------------------------------
%  Frode-bibelen -- 10 LesData.h -- verktøykassen
%  Inkluderes fra frode-bibelen.tex (kompiler alltid hovedfilen).
% ----------------------------------------------------------------------------

% ====================================================== 10 LESDATA.H ========
\kapittel{\texorpdfstring{\cc{LesData.h}}{LesData.h} -- verktøykassen}\label{kap:lesdata}

\cc{LesData.h} (EKS\_34/35, \href{https://frh.folk.ntnu.no/grprog/eksempler/LesData.h}{\cc{grprog/eksempler/LesData.h}})
inkluderes med \lstinline|#include "LesData.h"| og brukes i \emph{alle} løsningsforslag.
Den inkluderer selv \cc{ctype.h}, \cc{stdio.h}, \cc{stdlib.h} og \cc{string.h}.

\begin{center}
\begin{tabular}{@{}L{6.4cm}L{9.2cm}@{}}
\toprule
\textbf{Funksjon} & \textbf{Hva den gjør -- og typisk bruk} \\
\midrule
\cc{char lesChar(const char* t)} &
Skriver \cc{t:}, leser \emph{ett} tegn og returnerer det som \textbf{stor bokstav}. \newline
\cc{kommando = lesChar("\bs nValg");} \\
\cc{int lesInt(const char* t,} \newline \cc{\ \ \ \ \ \ \ \ const int min, const int max)} &
Skriver \cc{t (min - max):} og spør \textbf{igjen og igjen} til tallet er i $[\mathit{min}, \mathit{max}]$. \newline
\cc{nr = lesInt("\bs tKundenr", 1, gAntallKunder);} \\
\cc{float lesFloat(const char* t,} \newline \cc{\ \ \ \ \ \ \ \ const float min, const float max)} &
Som \cc{lesInt}, men for desimaltall (viser grensene med 2 desimaler). \newline
\cc{k->saldo = lesFloat("\bs tSaldo", 0, 100000);} \\
\cc{char* lagOgLesText(const char* t)} &
Leser en tekst, \textbf{\cc{malloc}-er nøyaktig nok plass} og returnerer pekeren. Må \cc{free}-es senere! \newline
\cc{k->navn = lagOgLesText("\bs tNavn");} \\
\cc{void lesText(const char* t,} \newline \cc{\ \ \ \ \ \ \ \ char* tekst, const int len)} &
Leser en tekst inn i et \textbf{ferdig buffer} (maks \cc{len}-1 tegn) og fjerner \cc{'\bs n'} på slutten. \newline
\cc{char buffer[STRLEN];  lesText("\bs tNavn", buffer, STRLEN);} \\
\bottomrule
\end{tabular}
\end{center}

\begin{felle}[{Feller med LesData.h}]
\begin{itemize}
  \item \cc{lesChar} returnerer \textbf{store} bokstaver: sammenlign med \cc{'J'}, \cc{'Q'}, \cc{'R'} -- aldri \cc{'j'}.
  \item \cc{lesInt(t, 1, gAntallKunder)} når \cc{gAntallKunder == 0} gir en \textbf{evig løkke} (ingen tall er mellom 1 og 0).
        Sjekk \lstinline|if (gAntallKunder > 0)| først -- eller bruk \cc{0} som «angre» (\cc{lesInt(t, 0, gAntallKunder)}).
  \item Brukeren skriver nummer \cc{1..antall}, arrayen bruker indeks \cc{0..antall-1}: \lstinline|gKunder[nr-1]|.
  \item \cc{lagOgLesText} lager \emph{ny} memory hver gang. Lagres teksten ikke (f.eks.\ duplikat), må du \cc{free} den.
  \item Alle funksjonene leser bort avsluttende \cc{'\bs n'} (med \cc{getchar()}), så de kan blandes fritt.
        Bruker du \cc{scanf} selv, må du også gjøre \lstinline|getchar();| før neste \cc{fgets}.
\end{itemize}
\end{felle}

\begin{tips}[{Ledetekster i Frodes stil}]
Ledetekster starter ofte med \cc{"\bs t"} for innrykk, og \cc{0} brukes som «ingen/angre»:\par
\lstinline|int nr = lesInt("\tFjern kunde (0 = angre)", 0, gAntallKunder);|
\end{tips}

\subsection{Kildekoden (din egen utskrift)}
\IfFileExists{LesData.h}{%
  \lstinputlisting[inputencoding=latin1,basicstyle=\ttfamily\scriptsize]{LesData.h}%
}{%
\begin{regel}[{Legg inn originalen her}]
Last ned \href{https://frh.folk.ntnu.no/grprog/eksempler/LesData.h}{\cc{LesData.h}} fra emnesiden,
legg filen i samme mappe som \cc{frode-bibelen.tex} og kompiler på nytt -- da settes hele kildekoden inn her automatisk.
(Alternativt: skriv den ut separat og legg den bak dette arket.)
\end{regel}
}
